% Book 3, Year 1, chapter 21 -- divide and conquer.
% Provenance (never printed): sources/book3/BRIEFS.md, chapter 21.
% ledger: b1:karatsuba, b1:strassen, b1:bfprt, b1:cpython-karatsuba
\chapter{Divide and Conquer}\label{ch:b1:divide-and-conquer}

Multiplying two numbers of 1000 digits by the method taught at school takes a million products of digits, one for
each digit of the first number and each digit of the second. For centuries nobody expected to do better: every digit
of one number must meet every digit of the other, it seems. It does not have to. Cutting each number in two halves and
combining three products of halves, instead of four, gives a method whose cost grows like $n^{1.585}$, and that is
how long numbers are multiplied in practice. The idea behind it, cutting a problem into smaller ones of the same kind,
has already given \omterm{def:b1:sorting:merge-sort}{merge sort}; this chapter turns it into a method and applies it to five problems.

\begin{recall}
The master theorem (\cref{thm:b1:asymptotic-analysis:master}) solves $T(n) = aT(n/b) + n^k$: $T(n) =
\Theta(n^k)$ if $a < b^k$, $\Theta(n^k \log n)$ if $a = b^k$, and $\Theta(n^{\log_b a})$ if $a > b^k$; a cost
$\Theta(n^k)$ in place of $n^k$ gives the same orders. \omterm{def:b1:sorting:merge-sort}{Merge sort} (\cref{def:b1:sorting:merge-sort}) satisfies $T(n) = 2T(n/2) + \Theta(n)$, hence $\Theta(n
\log n)$. An inversion of a sequence is a pair of positions $i < j$ whose keys are out of order
(\cref{def:b1:sorting:inversion}).
\end{recall}

\section{The method and its recurrences}\label{sec:b1:divide-and-conquer:method}

\begin{definition}[Divide and conquer]\label{def:b1:divide-and-conquer:dc}
A \emph{divide and conquer}\index{divide and conquer} algorithm solves an instance of a problem in three steps:
\emph{divide} it into smaller instances of the same problem, \emph{solve} each of them by a recursive call (or directly
when it is small enough), and \emph{combine} their solutions into a solution of the instance.
\end{definition}

The method was met informally in the second book, on searching a sorted list and sorting by merging. Its analysis is
always the same: if an instance of size $n$ is cut into $a$ instances of size $n/b$, and dividing plus combining cost
$f(n)$, the cost is $T(n) = aT(n/b) + f(n)$, and the master theorem gives its order. The work lies in finding a cheap
combine step, or in making $a$ small.

Counting inversions shows the three steps at their simplest. The inversions of a list are those inside its left half,
those inside its right half, and those with one key on each side. The recursive calls count the first two kinds; the
third kind is counted while merging the two sorted halves: when a key is taken from the right half, every key still
waiting in the left half is larger and comes before it.

\omcode{code/bachelor-1/21-divide-and-conquer/listings/inversions.py}{4}{19}{Counting inversions while
sorting}\label{lst:b1:divide-and-conquer:inversions}

\omoutput[2-2]{code/bachelor-1/21-divide-and-conquer/outputs/inversions.txt}{The sorted list and its 8
inversions.}\label{out:b1:divide-and-conquer:inversions}

The combine step costs $\Theta(n)$, so $T(n) = 2T(n/2) + \Theta(n) = \Theta(n \log n)$, against $\Theta(n^2)$ for
testing every pair. Sorting was not asked for, but it makes the combine step cheap: a \omterm{def:b1:divide-and-conquer:dc}{divide and conquer} algorithm
often returns more than the question, to make the next combination easy.

\section{Karatsuba's multiplication}\label{sec:b1:divide-and-conquer:karatsuba}

A number of $n$ digits is a list of digits, least significant first. The school method adds the product of digit $i$
of $a$ and digit $j$ of $b$ into cell $i + j$ of the result, then propagates the carries.

\omcodes{code/bachelor-1/21-divide-and-conquer/listings/digits.py}{16-33}{Carries, and the school method with its
counter of digit products}\label{lst:b1:divide-and-conquer:school}

Cut each number in two halves of $m = n/2$ digits: $a = a_1 10^m + a_0$ and $b = b_1 10^m + b_0$. Then
\[ ab = a_1 b_1 10^{2m} + (a_1 b_0 + a_0 b_1) 10^m + a_0 b_0, \]
four products of halves, and $T(n) = 4T(n/2) + \Theta(n)$ is $\Theta(n^2)$ again: cutting alone gains nothing. But
the middle term is a combination of the two others and of a single new product:
\[ a_1 b_0 + a_0 b_1 = a_1 b_1 + a_0 b_0 - (a_1 - a_0)(b_1 - b_0). \]

\begin{method}[Karatsuba's algorithm]\label{met:b1:divide-and-conquer:karatsuba}
To multiply two numbers of $n = 2m$ digits, compute recursively $z_2 = a_1 b_1$, $z_0 = a_0 b_0$ and $z_1 = |a_1 -
a_0| \cdot |b_1 - b_0|$, with $s$ the sign of $(a_1 - a_0)(b_1 - b_0)$, and return $z_2 10^{2m} + (z_2 + z_0 - s z_1)
10^m + z_0$. Numbers of one digit are multiplied directly.
\end{method}

The differences $|a_1 - a_0|$ and $|b_1 - b_0|$ are less than $10^m$: they have $m$ digits, like the halves, so that
the three recursive calls are on numbers of $m$ digits exactly. \index{Karatsuba's algorithm}

\omcodes{code/bachelor-1/21-divide-and-conquer/listings/karatsuba.py}{7-36}{The difference of two halves, and
Karatsuba's multiplication}\label{lst:b1:divide-and-conquer:karatsuba}

The combine loop adds and subtracts into cells that may go negative; \code{normalise} takes them back to digits with
floor division, and its \omterm{def:b1:c-modules-make-tests:assertion}{assertion} checks that the result fits in $2n$ digits, as the product of two $n$-digit numbers
must.

\begin{omfigure}
\begin{tikzpicture}[n/.style={draw, rounded corners=2pt, font=\scriptsize, inner sep=2pt, minimum width=7mm},
  level distance=12mm, level 1/.style={sibling distance=33mm}, level 2/.style={sibling distance=10mm}]
\node[n] (r) at (0,0) {$n$}
  child {node[n] {$n/2$} child {node[n] {$n/4$}} child {node[n] {$n/4$}} child {node[n] {$n/4$}}}
  child {node[n] {$n/2$} child {node[n] {$n/4$}} child {node[n] {$n/4$}} child {node[n] {$n/4$}}}
  child {node[n] {$n/2$} child {node[n] {$n/4$}} child {node[n] {$n/4$}} child {node[n] {$n/4$}}};
\node[font=\small] at (0,-3.4) {$\vdots$};
\node[font=\small] at (0,-4.0) {$3^k$ products of one digit};
\begin{scope}[font=\scriptsize, anchor=west]
\node at (5.0, 0) {1 call, $\Theta(n)$ to combine};
\node at (5.0,-1.2) {3 calls of $n/2$ \textcolor{omThm}{(4 with four products)}};
\node at (5.0,-2.4) {9 calls of $n/4$ \textcolor{omThm}{(16)}};
\node at (5.0,-4.0) {level $k = \log_2 n$ \textcolor{omThm}{($4^k = n^2$)}};
\end{scope}
\end{tikzpicture}
\par\smallskip{\small The \omterm{def:b1:asymptotic-analysis:recurrence}{recursion tree} of \omterm{met:b1:divide-and-conquer:karatsuba}{Karatsuba's algorithm} on $n = 2^k$ digits: three children per node. Level
$i$ holds $3^i$ calls on $n/2^i$ digits, and the combine work of the level is $(3/2)^i \Theta(n)$, growing to the
leaves. The same cut with four products (in red) has $4^i$ calls at level $i$.}
\end{omfigure}

The recurrence is now $T(n) = 3T(n/2) + \Theta(n)$, and $3 > 2^1$: the master theorem gives $\Theta(n^{\log_2 3})$,
with $\log_2 3 \approx 1.585$. The program counts the products exactly.

\omoutput{code/bachelor-1/21-divide-and-conquer/outputs/karatsuba.txt}{Digit products of the two methods, on a
small example and on $n$ nines times $n$ nines.}\label{out:b1:divide-and-conquer:karatsuba}

Fewer products do not mean less time: each call also pays for two differences, a combine loop and its carries, and a
call costs much more than a product of digits. Below some size, the school method wins. A \emph{cutoff} makes
Karatsuba's method fall back on the school method for numbers of at most that many digits.

\omoutput{code/bachelor-1/21-divide-and-conquer/outputs/karatsuba_timing.txt}{Microseconds per product of two random
$n$-digit numbers, on the author's laptop, for the school method and Karatsuba's with three
cutoffs.}\label{out:b1:divide-and-conquer:timing}

\begin{omfigure}
\begin{tikzpicture}
\begin{loglogaxis}[width=7.2cm, height=5.6cm, xlabel={digits $n$}, ylabel={digit products}, log basis x=2,
  legend pos=north west, legend style={font=\scriptsize}, tick label style={font=\scriptsize}]
\addplot[thick, omThm, mark=*, mark size=1.2pt] table[x=n, y=school]
  {code/bachelor-1/21-divide-and-conquer/figures/products.dat}; \addlegendentry{school, $n^2$}
\addplot[thick, omDef, mark=square*, mark size=1.2pt] table[x=n, y=karatsuba]
  {code/bachelor-1/21-divide-and-conquer/figures/products.dat}; \addlegendentry{Karatsuba, $n^{1.585}$}
\end{loglogaxis}
\end{tikzpicture}\hfill
\begin{tikzpicture}
\begin{loglogaxis}[width=7.2cm, height=5.6cm, xlabel={digits $n$}, ylabel={microseconds}, log basis x=2,
  legend pos=north west, legend style={font=\scriptsize}, tick label style={font=\scriptsize}]
\addplot[thick, omThm, mark=*, mark size=1.2pt] table[x=n, y=school]
  {code/bachelor-1/21-divide-and-conquer/figures/times.dat}; \addlegendentry{school}
\addplot[thick, black, dashed, mark=triangle*, mark size=1.2pt] table[x=n, y=c1]
  {code/bachelor-1/21-divide-and-conquer/figures/times.dat}; \addlegendentry{cutoff 1}
\addplot[thick, omDef, mark=square*, mark size=1.2pt] table[x=n, y=c32]
  {code/bachelor-1/21-divide-and-conquer/figures/times.dat}; \addlegendentry{cutoff 32}
\end{loglogaxis}
\end{tikzpicture}
\par\smallskip{\small Left: digit products counted in \cref{out:b1:divide-and-conquer:karatsuba}; the slopes are 2 and
$\log_2 3$. Right: the measured times of \cref{out:b1:divide-and-conquer:timing}; without a cutoff, Karatsuba's method
is still slower at 2048 digits, with a cutoff of 32 it is faster from 64 digits on.}
\end{omfigure}

\begin{history}[Three products instead of four]
The method was published by Karatsuba and Ofman in 1962. In 1969, Strassen applied the same idea to matrices: two $2
\times 2$ block matrices can be multiplied with 7 block products instead of 8, which gives $\Theta(n^{\log_2 7})$,
about $n^{2.81}$ operations, for $n \times n$ matrices.
\end{history}

\begin{dated}{2026}
CPython multiplies its integers by the school method unless both have more than 70 internal digits, of 30 bits each,
about 630 decimal digits; above that it uses Karatsuba's method. This is read in its source code, and such thresholds
change between versions.
\end{dated}

\section{The closest pair of points}\label{sec:b1:divide-and-conquer:closest}

\begin{definition}[Closest-pair problem]\label{def:b1:divide-and-conquer:closest}
The \emph{closest-pair problem}\index{closest-pair problem} asks, for $n \geq 2$ distinct points of the plane, for the
smallest distance between two of them.
\end{definition}

Comparing every pair takes $n(n - 1)/2$ distances. \omterm{def:b1:divide-and-conquer:dc}{Divide and conquer} sorts the points by $x$, cuts them by a vertical
line through the median $x$, and solves both halves: $\delta$ is the smaller of the two answers. A closer pair can only
have one point on each side, both within $\delta$ of the line: in the \emph{strip} of width $2\delta$ around it.

\begin{proposition}[Seven neighbours suffice]\label{prop:b1:divide-and-conquer:seven}
Sort the points of the strip by $y$. If two of them, $p$ before $q$, are less than $\delta$ apart, then $q$ is one of
the seven points that follow $p$.
\end{proposition}

\begin{proof}
All the points that follow $p$ and are within $\delta$ of it lie in the rectangle of width $2\delta$ and height $\delta$
whose bottom edge passes through $p$, centred on the line. Cut this rectangle into eight squares of side $\delta/2$,
four on each side of the line. Two points in one square are at most $\delta/\sqrt{2} < \delta$ apart, and on one side
of the line, which contradicts the definition of $\delta$: each square holds at most one point. With $p$ itself, the
rectangle holds at most eight points, so at most seven points come after $p$ before the first one with $y \geq p_y +
\delta$.
\end{proof}

\begin{omfigure}
\begin{tikzpicture}[scale=1.1]
\fill[omDef!10] (-2,0) rectangle (2,2);
\draw[dashed] (0,-0.6) -- (0,2.7) node[above, font=\scriptsize] {dividing line};
\draw[thin, gray] (-2,-0.6) -- (-2,2.7) (2,-0.6) -- (2,2.7);
\foreach \i in {-2,-1,0,1} \foreach \j in {0,1} \draw[omDef] (\i,\j) rectangle (\i+1,\j+1);
\fill[omThm] (-0.6,0) circle (2pt) node[below, font=\scriptsize] {$p$};
\fill (0.2,0.3) circle (1.5pt) (1.9,2) circle (1.5pt) (-1.9,1.9) circle (1.5pt);
\draw[<->] (-2,-0.45) -- (0,-0.45) node[midway, below, font=\scriptsize] {$\delta$};
\draw[<->] (2.2,0) -- (2.2,2) node[midway, right, font=\scriptsize] {$\delta$};
\node[font=\scriptsize, anchor=west] at (2.6,1.4) {eight squares of side $\delta/2$:};
\node[font=\scriptsize, anchor=west] at (2.6,1.0) {at most one point in each};
\end{tikzpicture}
\par\smallskip{\small The strip argument. The points that follow $p$ in the $y$ order and are within $\delta$ of it lie
in the shaded $2\delta \times \delta$ rectangle; each of its eight squares holds at most one point. Points on one side
are at least $\delta$ apart; the point just right of the line is the one that can beat $\delta$.}
\end{omfigure}

\omcode{code/bachelor-1/21-divide-and-conquer/listings/closest_pair.py}{27}{45}{The closest pair, returning the points
sorted by $y$ for the caller's strip}\label{lst:b1:divide-and-conquer:closest}

Each call returns its points sorted by $y$, merged from the two halves, so that the strip is sorted without sorting:
the combine step costs $\Theta(n)$ and the whole $\Theta(n \log n)$. On the town's 30 places and on random points:

\omoutput{code/bachelor-1/21-divide-and-conquer/outputs/closest_pair.txt}{The closest pair of the town's places, and
distances computed on random points, against \omterm{def:b1:backtracking:exhaustive}{brute force}.}\label{out:b1:divide-and-conquer:closest}

On random points the strip is thin and holds few points: barely more than one distance per point.

\section{Selection in linear time}\label{sec:b1:divide-and-conquer:selection}

\begin{definition}[Selection problem, quickselect, median of medians]\label{def:b1:divide-and-conquer:selection}
The \emph{selection problem}\index{selection problem} asks for the key of rank $k$ of $n$ keys, the one that would be
at position $k$ if they were sorted; the median is the key of rank $\lfloor (n - 1)/2 \rfloor$. \emph{Quickselect}
\index{quickselect} partitions the keys around a pivot, as \omterm{def:b1:sorting:quicksort}{quicksort} does, and goes on in the one part that holds rank
$k$. The \emph{median of medians}\index{median of medians} chooses as pivot the median of the medians of groups of five
keys, found by a recursive call.
\end{definition}

Sorting answers in $\Theta(n \log n)$; selection needs less. \omterm{def:b1:divide-and-conquer:selection}{Quickselect} makes one recursive call instead of two, so
with pivots that halve the keys its cost is $n + n/2 + n/4 + \dots \leq 2n$ partition steps. With bad pivots it is
quadratic, like \omterm{def:b1:sorting:quicksort}{quicksort}.

\omcodes{code/bachelor-1/21-divide-and-conquer/listings/selection.py}{15-31,45-57}{Partition, \omterm{def:b1:divide-and-conquer:selection}{quickselect}, and the
\omterm{def:b1:divide-and-conquer:selection}{median of medians}}\label{lst:b1:divide-and-conquer:select}

\omoutput{code/bachelor-1/21-divide-and-conquer/outputs/selection.txt}{Comparisons to find the key of rank $n/2$ (the
upper median) of $n$ keys, sorted or shuffled.}\label{out:b1:divide-and-conquer:select}

On a sorted list, the last key as pivot is the largest, and \omterm{def:b1:divide-and-conquer:selection}{quickselect} removes one key per round: rounds on $n$,
$n - 1$, \dots, $n/2$ keys, about $3n^2/8$ keys partitioned in all, two comparisons per key: about $3n^2/4$, the 750\,500 of the
table for $n = 1000$. A random pivot gives a linear cost on average. The \omterm{def:b1:divide-and-conquer:selection}{median of
medians} gives it in the worst case.

\begin{omfigure}
\begin{tikzpicture}[scale=0.55]
\fill[omDef!20] (-0.4,-0.4) rectangle (3.4,2.4);
\fill[omThm!15] (2.6,1.6) rectangle (6.4,4.4);
\foreach \c in {0,...,6} \foreach \r in {0,...,4} \fill (\c,\r) circle (3pt);
\foreach \c in {0,...,6} \draw[gray] (\c,0) -- (\c,4);
\draw[very thick, omThm] (3,2) circle (8pt);
\draw[dashed] (-0.8,2) -- (6.8,2);
\node[font=\scriptsize, anchor=west] at (7.2,2) {medians of the columns, in increasing order};
\node[font=\scriptsize, anchor=west] at (7.2,4) {each column sorted, smallest at the bottom};
\node[font=\scriptsize, anchor=west, omDef] at (7.2,0.5) {at most the pivot: $3 \cdot 4 = 12$ keys here};
\node[font=\scriptsize, anchor=west, omThm] at (7.2,3.2) {at least the pivot: 12 keys};
\end{tikzpicture}
\par\smallskip{\small Thirty-five keys in seven columns of five, columns ordered by their median; the circled key, the
median of the medians, is the pivot. The shaded block below and left holds keys no larger than the pivot, the one
above and right keys no smaller: each about $3/10$ of all keys.}
\end{omfigure}

\begin{theorem}[Selection in linear time]\label{thm:b1:divide-and-conquer:mom}
The \omterm{def:b1:divide-and-conquer:selection}{median of medians} finds the key of rank $k$ of $n$ keys in $O(n)$ comparisons.
\end{theorem}

\begin{proof}
Of the $g = \lceil n/5 \rceil$ medians, at least $\lceil g/2 \rceil$ are at most the pivot, and each of their groups,
except perhaps the pivot's own and a last incomplete one, holds three keys at most the pivot. So at least $3(\lceil
g/2 \rceil - 2) \geq 3n/10 - 6$ keys are at most the pivot; likewise at least that many are at least the pivot. The
recursive call on one side therefore receives at most $7n/10 + 6$ keys. Sorting the groups and partitioning cost at
most $an$ comparisons for a constant $a$, so
\[ T(n) \leq T(\lceil n/5 \rceil) + T(7n/10 + 6) + an. \]
Suppose $T(j) \leq cj$ for $j < n$, with $n \geq 140$ and $c = 20a$, and $T$ bounded by $cj$ on small sizes by taking
$c$ larger if needed. Then $T(n) \leq c(n/5 + 1) + c(7n/10 + 6) + an = cn - c(n/10 - 7) + an$, and $n/10 - 7 \geq
n/20$ for $n \geq 140$, so $T(n) \leq cn - cn/20 + an = cn$.
\end{proof}

The proof works because $1/5 + 7/10 < 1$: the two calls together receive less than the whole. Groups of three fail
this test (\cref{exo:b1:divide-and-conquer:10}). The constant is large: eight to eleven comparisons per key in
\cref{out:b1:divide-and-conquer:select}, against three to seven for \omterm{def:b1:divide-and-conquer:selection}{quickselect} on shuffled keys or with random
pivots. A practical routine runs \omterm{def:b1:divide-and-conquer:selection}{quickselect} and can fall back on a guaranteed method when the recursion goes too deep,
as \omterm{def:b1:sorting:introsort}{introsort} does for sorting. The algorithm is due to Blum, Floyd, Pratt, Rivest and Tarjan, in 1973.

\section{Dichotomy on the answer, and meeting in the middle}\label{sec:b1:divide-and-conquer:dichotomy}

Binary search halves a sorted \omterm{def:b1:pointers-and-arrays:array}{array} (\cref{prop:b1:program-correctness:bsearch}). It halves anything ordered, including
the answers to a question.

\begin{definition}[Monotone predicate, dichotomy on the answer]\label{def:b1:divide-and-conquer:dichotomy}
A predicate $P$ on the integers is \emph{monotone}\index{monotone predicate} if $P(x)$ implies $P(x + 1)$: false up to
some point, true after it. A \emph{dichotomy on the answer}\index{dichotomy on the answer} finds the smallest $x$ in
$[lo, hi]$ with $P(x)$, knowing $P(hi)$, by binary search on $x$, with $\lceil \log_2 (hi - lo + 1) \rceil$ evaluations
of $P$.
\end{definition}

The town wants to cover its 30 places with $k$ snow-clearing zones, each a band of the same width from west to east.
For a given width, a greedy pass decides whether $k$ bands suffice: start a band at the westernmost place not yet
covered. A larger width never needs more bands, so the predicate is \omterm{def:b1:divide-and-conquer:dichotomy}{monotone}, and the smallest width is found by
dichotomy.

\omcodes{code/bachelor-1/21-divide-and-conquer/listings/segments.py}{10-31}{The greedy test and the dichotomy on the
width}\label{lst:b1:divide-and-conquer:segments}

\omoutput{code/bachelor-1/21-divide-and-conquer/outputs/segments.txt}{The smallest width of $k$ bands covering the
town's places, and the number of tests.}\label{out:b1:divide-and-conquer:segments}

Twelve or thirteen tests decide among 5625 widths. The method applies whenever checking an answer is easier than
finding it, and the check is \omterm{def:b1:divide-and-conquer:dichotomy}{monotone}.

\begin{definition}[Meet in the middle]\label{def:b1:divide-and-conquer:meet}
\emph{Meeting in the middle}\index{meet in the middle} splits a search over choices into two halves, lists the partial
results of each half, and finds the compatible pairs by sorting or hashing, instead of trying every combination.
\end{definition}

The subset-sum search of \cref{ch:b1:backtracking} visits up to $2^n$ subsets. Split the $n$ values into two halves:
each has $2^{n/2}$ subset sums; a subset of all the values sums to the target exactly when its left part sums to $s$ and
its right part to $\mathit{target} - s$. A table of how many right subsets give each sum answers that question for all
$s$.

\omcodes{code/bachelor-1/21-divide-and-conquer/listings/meet_middle.py}{16-30}{Subset sum by meeting in the
middle}\label{lst:b1:divide-and-conquer:meet}

\omoutput{code/bachelor-1/21-divide-and-conquer/outputs/meet_middle.txt}{Reading lists of exactly 1000 pages from the
first $n$ books.}\label{out:b1:divide-and-conquer:meet}

The counts for 16 and 24 books are those found by \omterm{def:b1:backtracking:backtracking}{backtracking} in \cref{out:b1:backtracking:subset-sum}. The cost is
$\Theta(2^{n/2})$ sums and table operations, with a table of the same size: memory is traded for time. Forty books are
reached with two million sums, where the full tree has a trillion leaves.

\section{Exercises}\label{sec:b1:divide-and-conquer:exercises}

\begin{exercise}[$\star$]\label{exo:b1:divide-and-conquer:1}
List the inversions of $[3, 1, 4, 1, 5, 9, 2, 6]$ and check the count of \cref{out:b1:divide-and-conquer:inversions}.
How many are counted by the top-level merge?
\end{exercise}

\begin{exercise}[$\star$]\label{exo:b1:divide-and-conquer:2}
Give the three products of halves that Karatsuba's method computes at the top level for $1234 \times 5678$, the sign
$s$, and the result.
\end{exercise}

\begin{exercise}[$\star$]\label{exo:b1:divide-and-conquer:3}
Run \code{quickselect} by hand on $[7, 2, 9, 4, 1, 8, 5]$ with $k = 5$ and the last key as pivot: give the pivots and
the parts kept.
\end{exercise}

\begin{exercise}[$\star$]\label{exo:b1:divide-and-conquer:4}
For the points $0, 3, 4, 10, 12, 20$ and $k = 2$, give the values of \code{covers} for every width, and the widths the
dichotomy tests.
\end{exercise}

\begin{exercise}[$\star\star$]\label{exo:b1:divide-and-conquer:5}
This version compares each point of the strip with the next one only. Explain the counterexample it finds.
\omcode{code/bachelor-1/21-divide-and-conquer/listings/exo_strip_bug.py}{7}{16}{A strip checked too lightly}
\omoutput[2-3]{code/bachelor-1/21-divide-and-conquer/outputs/exo_strip_bug.txt}{Against \omterm{def:b1:backtracking:exhaustive}{brute force}.}
\end{exercise}

\begin{exercise}[$\star\star$]\label{exo:b1:divide-and-conquer:6}
This dichotomy computes a cube root with an absolute tolerance. Explain why it never ends for $y = 10^{15}$, and repair
it.
\omcode{code/bachelor-1/21-divide-and-conquer/listings/exo_float_bisect.py}{4}{18}{A dichotomy on reals}
\omoutput[2-4]{code/bachelor-1/21-divide-and-conquer/outputs/exo_float_bisect.txt}{Three cube roots.}
\end{exercise}

\begin{exercise}[$\star\star$]\label{exo:b1:divide-and-conquer:7}
Compute the largest sum of a non-empty slice of a list by \omterm{def:b1:divide-and-conquer:dc}{divide and conquer}, in $O(n \log n)$, and test it against a
linear method that keeps the best sum ending at each position.
\end{exercise}

\begin{exercise}[$\star\star$]\label{exo:b1:divide-and-conquer:8}
A \emph{majority key} of a list is held by more than half of its positions. Find it, or report that there is none, by
\omterm{def:b1:divide-and-conquer:dc}{divide and conquer} in $O(n \log n)$.
\end{exercise}

\begin{exercise}[$\star\star$]\label{exo:b1:divide-and-conquer:9}
Karatsuba's method pads numbers of $n$ digits to the next power of 2. Show that it makes fewer than $3 n^{\log_2 3}$
digit products, and that with a cutoff $2^j$ it makes $3^{k-j} 4^j$ on $2^k$ digits.
\end{exercise}

\begin{exercise}[$\star\star\star$]\label{exo:b1:divide-and-conquer:10}
Show that the \omterm{def:b1:divide-and-conquer:selection}{median of medians} with groups of three gives the recurrence $T(n) \leq T(n/3) + T(2n/3) + an$, and that
this recurrence does not imply $T(n) = O(n)$: its solution with equality is $\Theta(n \log n)$.
\end{exercise}

\begin{exercise}[$\star\star\star$]\label{exo:b1:divide-and-conquer:11}
Buildings are rectangles $(\mathit{left}, \mathit{right}, \mathit{height})$ standing on a line. Compute the skyline,
the list of points where the visible height changes, by \omterm{def:b1:divide-and-conquer:dc}{divide and conquer} in $O(n \log n)$.
\end{exercise}

\begin{exercise}[$\star\star\star$]\label{exo:b1:divide-and-conquer:12}
Find four positions $i < j < k < m$ of a list whose keys sum to a target in $O(n^2 \log n)$, by meeting in the middle
on the sums of pairs, and prove the method correct.
\end{exercise}

\section{Problem: Multiplying long numbers}\label{sec:b1:divide-and-conquer:problem}

\begin{problem}[{Weekend problem --- Multiplying long numbers}]\label{pb:b1:divide-and-conquer:1}
The listings of this chapter multiply numbers given as lists of decimal digits. The problem counts their work exactly,
then reads their measured times (\cref{out:b1:divide-and-conquer:karatsuba,out:b1:divide-and-conquer:timing}).

\needspace{4\baselineskip}
\medskip\noindent\textbf{Part I --- The school method.}
\begin{enumerate}
\item How many digit products does the school method make for two numbers of $n$ digits? Check with the output.
\item Show that a cell of \code{r}, before the carries, holds at most $81n$.
\item Show that propagating the carries once, at the end, costs $\Theta(n)$ and gives the right digits.
\item Why does the result of two $n$-digit numbers fit in $2n$ digits?
\item By what factor does the school method's time grow when $n$ doubles, according to the table? What does the
model predict?
\item How many digit products does the opening's multiplication of two 1000-digit numbers take?
\end{enumerate}

\needspace{4\baselineskip}
\medskip\noindent\textbf{Part II --- Karatsuba's identity.}
\begin{enumerate}[resume]
\item Expand $ab$ with $a = a_1 10^m + a_0$, $b = b_1 10^m + b_0$.
\item Prove $a_1 b_0 + a_0 b_1 = a_1 b_1 + a_0 b_0 - (a_1 - a_0)(b_1 - b_0)$.
\item Why do $|a_1 - a_0|$ and $|b_1 - b_0|$ have at most $m$ digits?
\item What does \code{difference} return, and why does it compare from the top digit?
\item The combine loop leaves cells that may be negative. Why does \code{normalise} still give the right digits?
\item Write the recurrence for the number $K(n)$ of digit products, $n$ a power of 2.
\end{enumerate}

\needspace{4\baselineskip}
\medskip\noindent\textbf{Part III --- Exact counts.}
\begin{enumerate}[resume]
\item Prove $K(2^k) = 3^k$.
\item Show that $3^k = n^{\log_2 3}$ for $n = 2^k$.
\item Compare $K(1024)$ with $1024^2$.
\item Show that the other work of one call on $n$ digits is $\Theta(n)$, and that the total work is still
$\Theta(n^{\log_2 3})$.
\item With a cutoff $c = 2^j$, how many digit products are made on $2^k$ digits (\cref{exo:b1:divide-and-conquer:9})?
Give the count for $n = 1024$ and $c = 32$.
\item Which cutoff minimises the number of digit products? Why is a larger one used anyway?
\end{enumerate}

\needspace{4\baselineskip}
\medskip\noindent\textbf{Part IV --- Where it pays, on this machine.}
\begin{enumerate}[resume]
\item Without a cutoff, by what factor is Karatsuba's method slower than the school method at each $n$ of the table?
How does that factor change when $n$ doubles?
\item Estimate from that trend the size from which Karatsuba's method without a cutoff would win.
\item From which $n$ of the table does Karatsuba's method with cutoff 8, then with cutoff 32, beat the school method?
\item By what factor does the time with cutoff 32 grow when $n$ doubles at the end of the table? Compare with $3$ and
$4$.
\item CPython's cutoff is 70 digits of 30 bits. How many decimal digits is that?
\item State the named result, \emph{Karatsuba's count}.
\end{enumerate}
\end{problem}
